\section{Related Work}

\paragraph{Video behavioral recognition.} DEPART studies depression and Parkinson detection from in-the-wild video with body-centric CLIP features, temporal modeling, and prototype-aware prediction \citep{ryumina2026depart}. Its results motivate our two bounded video families and our attention to detection and modality failures. However, the published DEPART formulation combines data into mutually exclusive classes, whereas WSM keeps two independent masked disease heads. This distinction is central to the comparison boundary in Section~\ref{sec:results}.

\paragraph{Multimodal learning with disjoint corpora.} Task-aware cross-modal experts such as TACME motivate directed multimodal representations and task-conditioned aggregation \citep{ryumina2025tacme}. In WSM, corpus identity is structurally coupled to the observed task, so task-aware mechanisms and corpus diagnostics are treated as hypotheses to be ablated rather than assumptions.

\paragraph{Reliability and multi-objective learning.} Smooth Tchebycheff scalarization is a computationally economical approach to balancing competing objectives \citep{lin2024stch}; gradient-based multi-objective learning offers a broader taxonomy \citep{chen2025moo}. Our evidence does not establish a general advantage for the tested RA-STCH controller over simpler balancing. It does support the narrower reliability contribution under the frozen ablation gate. Transcript context uses a frozen XLM-R baseline \citep{conneau2020xlmr}; visual features follow the frozen CLIP feature family \citep{radford2021clip}.
